\section{Background}
\label{sec:background}

\subsection{Messages and derivations}
Messages $\Msgs$ are terms of the free algebra over atoms (names, fresh values, public constants), pairing $\pair{a}{b}$, symmetric encryption $\senc{m}{k}$, asymmetric encryption $\aenc{m}{\pk{s}}$, and public and private function symbols. The inference rules $\RR$ are the usual Dolev--Yao rules~\cite{dolev1983security}: the constructors $\lbl{pair}$, $\lbl{senc}$, $\lbl{aenc}$ and one constructor for each public function symbol, and the destructors $\lbl{proj}$ (both projections), $\lbl{sdec}$ (premises: ciphertext and key) and $\lbl{adec}$ (premises: ciphertext and private key). The first premise of a destructor is its \emph{principal} premise.

A \emph{derivation} $\pi$ of a target message $t$ from a finite set $K\subseteq\Msgs$ is a finite tree whose leaves are elements of $K$, whose inner nodes apply rules of $\RR$, and whose root is $t$. Its \emph{hypotheses} $\Hyp_\pi\subseteq K$ are its leaves and $\Rules(\pi)\subseteq\RR$ is the set of rules it applies. A derivation is \emph{normal} if no destructor's principal premise is the conclusion of the inverse constructor and no constructor rebuilds a term from the outputs of its inverse destructor. $\Norm_t(K)$ is the set of normal derivations of $t$ from $K$.

\begin{lemma}[Subterm property]\label{lem:subterm}
Every message in a derivation $\pi\in\Norm_t(K)$ belongs to $\mathrm{St}(K\cup\{t\})$, where \(\mathrm{St}(X)\) is the set of subterms of the terms in \(X\).
\end{lemma}
\begin{proof}
Let $v$ be the principal premise of a destructor. The message at $v$ has top symbol $\langle\cdot,\cdot\rangle$, $\mathsf{senc}$ or $\mathsf{aenc}$, so in the free algebra the only constructor that can produce it is the inverse one, which normality excludes. Hence $v$ is a leaf or the conclusion of a destructor, and in the latter case the same argument applies to that destructor's principal premise. Following principal premises from $v$ we reach a leaf $x\in K$; since the conclusion of a destructor is a subterm of its principal premise, the message at $v$ is in $\mathrm{St}(x)\subseteq\mathrm{St}(K)$.

We show by induction on the distance from the root that the message at every node $v$ is in $\mathrm{St}(K\cup\{t\})$. The root carries $t$. Otherwise let $p$ be the parent of $v$. If $p$ is a constructor, the message at $v$ is an argument of the message at $p$, which is in $\mathrm{St}(K\cup\{t\})$ by induction. If $v$ is the principal premise of $p$, the first part applies. If $v$ is the key premise of $p$, its message is the key $k$ of $\senc{m}{k}$ or the private key $s$ of $\mathsf{aenc}(m,\mathsf{pk}(s))$, a subterm of the principal premise, which is in $\mathrm{St}(K)$ by the first part.
\end{proof}

\begin{remark}[Finite basis]\label{rem:finite}
$\Norm_t(K)$ can be infinite: for $K=\{\senc{a}{b},\senc{b}{a},a\}$, $a$ is obtained by decrypting $\senc{a}{b}$ with $b$, which is obtained by decrypting $\senc{b}{a}$ with $a$, and so on. We call $\Basis^*_t(K)\subseteq\Norm_t(K)$, the normal derivations in which no message occurs twice on a root-to-leaf path, the \emph{finite basis}. It is finite. By \Cref{lem:subterm}, every message in a normal derivation lies in $\mathrm{St}(K\cup\{t\})$, which is finite since $K$ is. Hence a root-to-leaf path of a derivation in $\Basis^*_t(K)$ has at most $|\mathrm{St}(K\cup\{t\})|$ nodes. Every node has at most $d$ children, where $d$ is the larger of $2$ and the maximal arity of a function symbol in $\mathrm{St}(K\cup\{t\})$, and is labelled by a message of $\mathrm{St}(K\cup\{t\})$. Trees of bounded height and bounded branching with labels from a finite set are finitely many, so $\Basis^*_t(K)$ is finite.
\end{remark}

\subsection{The Forget mutation}\label{subsec:forget}
In a security ceremony, we say that the role of the human can be represented by a sequence of transitions $\AgSt(\human,i,\K{\human}{i}),\ \mathit{Rcv}(m_1) \xrightarrow{a} \AgSt(\human,i+1,\K{\human}{i+1}),\ \mathit{Snd}(m_2)$, where we define $\K{\human}{i}$ as the human's knowledge (monotone), and where $a$ contains internal actions. \forget{} records forgotten messages in a set $\Forget^i_\human$. Whether a hypothesis $\bar m$ is \emph{blocked} depends on a closure policy $\kappa$: under Case~1, $\bar m\in\Forget$; under Case~2, $\bar m$ is derivable from $\Forget$ with the pairing rules; under Case~3, with all Dolev--Yao rules. $\Blocked(X,\Forget)$ is the set of blocked elements of $X$. Blocking is pointwise: whether $\bar m$ is blocked depends on $\bar m$, $\Forget$ and $\kappa$ only. A message $t$ is \emph{derivable under forgetting} if some derivation $\pi$ of $t$ has $\Blocked(\Hyp_\pi,\Forget)=\varnothing$~\cite[Def.~4]{sandu2026formal}.

Forgetting, in the set-based algorithm, works by transitions: a forgotten known message is added to $\Forget$; a forgotten received $m_1$ is not added to the knowledge. If the forgotten message occurs in an internal action, the action is removed (\neglect). If some derivation $m_2$ has no blocked hypothesis, $m_2$ is sent; if not, for each derivation and blocked hypothesis $\bar m$ it collects the set $R_{\bar m}$ of known messages with the same format and type label, and emits one \replace{} per combination of replacements, or \skipm{} if there is none. Re-receiving a forgotten message, or deriving it from a received one, removes it from $\Forget$.

\subsection{Defeasible Logic}
A defeasible theory $\Dth=(F,R,>)$ has a finite set $F$ of literals (facts), a finite set $R$ of labelled rules $r: A(r)\hookrightarrow C(r)$ that are strict ($\strict$), defeasible ($\defeas$) or defeaters ($\defeat$), and an acyclic superiority relation $>$ on rules with complementary heads~\cite{nute1994defeasible,antoniou2001representation}. A defeater cannot support its head and can only prevent the complementary conclusion. We write $R[q]$ for the rules with head $q$, $R_s[q]$ and $\Rsd[q]$ for its strict, and strict or defeasible, rules, and ${\sim}q$ for the complement of $q$.
We use the ambiguity-blocking logic DL($\partial$) with \emph{team defeat}~\cite{antoniou2001representation}. Its conclusions are tagged: $\pD q$ and $\mD q$ for definite (strict) provability and its refutation, and
\begin{align*}
\pd q \iff{}& \pD q \text{ or } \bigl[\exists r\in\Rsd[q]\ \forall a\in A(r){:}\ \pd a\bigr] \wedge \mD{\sim}q \wedge \forall s\in R[{\sim}q]{:}\\
&\quad \bigl(\exists a\in A(s){:}\ \md a\bigr) \vee \bigl(\exists t\in\Rsd[q]{:}\ \forall a\in A(t){:}\ \pd a \wedge t>s\bigr),
\end{align*}
with $\md q$ defined as its strong negation (see \cref{app:variants} for $\md q$, $\pD q$ and $\mD q$). A rule is \emph{applicable} if all its antecedents are $\pd$, and \emph{discarded} if one is $\md$. Two features matter below. \emph{Team defeat}: each attacker $s$ of $q$ may be overruled by a \emph{different} applicable supporter $t$. \emph{Ambiguity blocking}: a literal whose support is tied with its attack receives $\md$, and rules depending on it are discarded.
